\documentclass[a4paper,12pt]{article}
\usepackage[T2A]{fontenc}
\usepackage[utf8]{inputenc}
\usepackage[russian]{babel}
\usepackage{amsmath,amssymb,mathtools}
\usepackage{geometry}
\geometry{top=22mm,bottom=22mm,left=24mm,right=22mm}
\renewcommand{\familydefault}{\sfdefault}
\usepackage{microtype}
\usepackage{array,booktabs,tabularx}
\usepackage{enumitem}
\usepackage{parskip}
\usepackage{xcolor}
\usepackage{hyperref}
\usepackage{tikz}
\usetikzlibrary{calc}
\definecolor{accent}{HTML}{007A7A}
\definecolor{deep}{HTML}{134F50}
\definecolor{textgray}{HTML}{343D43}
\definecolor{errorred}{HTML}{A02C30}
\hypersetup{colorlinks=true,linkcolor=deep,urlcolor=deep,pdftitle={Ревью домашней работы — 13 ноября 2026 года},pdfauthor={Лёвин Артём Александрович}}
\setlength{\parindent}{0pt}
\setlength{\parskip}{6pt}
\setlist[itemize]{leftmargin=*,nosep}
\newcommand{\topic}[1]{\par\vspace{5pt}{\large\bfseries\color{deep}#1}\par\vspace{2pt}}
\newcommand{\err}{\textcolor{errorred}{\textbf{Ошибка:}}}
\newcommand{\keyidea}{\textcolor{blue!60!black}{\textbf{Ключевая идея:}}}
\newcommand{\finalword}{\textcolor{accent}{\textbf{Итог:}}}
\newcommand{\vp}[2]{\nu_{#1}(#2)}
\newcommand{\floor}[1]{\left\lfloor #1\right\rfloor}
\pagestyle{plain}
\begin{document}
\thispagestyle{empty}
\begin{center}
\vspace*{47mm}
{\Large\color{deep} МАТЕМАТИКА}\par\vspace{18mm}
{\Huge\bfseries Ревью домашней работы}\par\vspace{9mm}
{\Large Факториалы и делимость степеней}\par\vspace{26mm}
{\large Автор: Лёвин Артём Александрович}\par\vspace{6mm}
{\large 13 ноября 2026 года}
\end{center}
\vfill
\begin{center}
\textcolor{textgray}{\small Лёвин Артём Александрович эксклюзивно для Ксении Клыковой}\par\vspace{8mm}
\begin{tikzpicture}[x=1cm,y=1cm]
\draw[accent!65,line width=0.9pt] (0,0) .. controls (3.0,0.55) and (7.0,-0.55) .. (10.5,0);
\end{tikzpicture}
\end{center}
\newpage
\begin{center}
{\Large\bfseries\color{deep} Сводная таблица ошибок}
\end{center}
\vspace{8mm}
\small
\renewcommand{\arraystretch}{1.28}
\begin{tabularx}{\linewidth}{@{}p{15mm}p{24mm}X p{18mm}p{19mm}@{}}
\toprule
\textbf{№} & \textbf{Тема} & \textbf{Суть неточности} & \textbf{Ваш ответ} & \textbf{Верный ответ}\\
\midrule
\hyperref[task:8]{8} & Факториал & В записи для $15!$ пропущены простые множители $11$ и $13$. & $5$ & $5$\\
\addlinespace[4pt]
\hyperref[task:9]{9} & Степень делителя & В записи $18^k$ потерян показатель $k$; разложение $20!$ неполное. & $4$ & $4$\\
\addlinespace[4pt]
\hyperref[task:10]{10} & Степени простых & В записи $72^k$ потерян показатель $k$; неверно найдено $\nu_5(30!)$. & $7$ & $7$\\
\addlinespace[4pt]
\bottomrule
\end{tabularx}
\par\normalsize
\vspace{13mm}
Во всех трёх заданиях числовой ответ получен верно. Замечания касаются математической точности промежуточных равенств. При использовании формулы делимости достаточно вычислять показатели простых чисел, входящих в основание степени; полное разложение факториала записывать необязательно.
\newpage
\thispagestyle{empty}
\vspace*{0.34\textheight}
\begin{center}
{\Huge\bfseries\color{deep} Разбор ошибок}
\end{center}
\newpage
\phantomsection\label{task:8}
\begin{center}{\LARGE\bfseries\color{deep} Задание № 8}\end{center}
\topic{Текст задания}
Найдите наибольшее целое неотрицательное $k$, при котором $12^k\mid 15!$.
\topic{Ваше решение}
Верно представлены $12^k=2^{2k}3^k$ и показатели простых чисел:
\[
\vp{2}{15!}=7+3+1=11,\qquad \vp{3}{15!}=5+1=6.
\]
Также вычислены $\vp{5}{15!}=3$ и $\vp{7}{15!}=2$. Затем записано
\[
15!=2^{11}\cdot3^6\cdot5^3\cdot7^2,
\]
после чего найдено $k=\min\{\floor{11/2},\floor{6/1}\}=5$.
\topic{Ваш ответ}
$k=5$ — верно.
\topic{Ошибка}
\err{} Неполное разложение факториала записано как точное равенство.
\topic{Где именно возникает ошибка}
\textbf{Последний правильный шаг:} показатели при двойке, тройке, пятёрке и семёрке в $15!$ вычислены правильно.

\textbf{Первый неправильный шаг:} равенство $15!=2^{11}\cdot3^6\cdot5^3\cdot7^2$. Число $15!$ дополнительно делится на $11$ и на $13$.
\topic{Почему это неверно}
По определению $15!=1\cdot2\cdot\ldots\cdot15$. Среди этих множителей находятся простые числа $11$ и $13$, поэтому при \emph{полном} разложении их опустить нельзя. Записанное произведение представляет только часть простых множителей и меньше самого факториала. При этом для решения исходной задачи показатели при $11$ и $13$ действительно не требуются: основание $12=2^2\cdot3$ содержит только двойки и тройки.
\topic{Ключевая идея}
\keyidea{} Следует различать полное равенство для факториала и выписанные отдельно показатели нужных простых чисел $\vp{p}{n!}$.
\topic{Исправление}
Если требуется полное разложение, нужно дописать недостающие множители:
\[
\boxed{15!=2^{11}\cdot3^6\cdot5^3\cdot7^2\cdot11\cdot13.}
\]
В решении задачи можно вообще обойтись без этого равенства.
\topic{Полное правильное решение}
Для делимости $12^k\mid15!$ представим основание степени:
\[
12^k=(2^2\cdot3)^k=2^{2k}\cdot3^k.
\]
Показатель двойки в факториале вычисляется по формуле Лежандра:
\[
\vp{2}{15!}=\floor{15/2}+\floor{15/4}+\floor{15/8}
=7+3+1=11.
\]
Показатель тройки:
\[
\vp{3}{15!}=\floor{15/3}+\floor{15/9}=5+1=6.
\]
Следовательно, для делимости одновременно нужны два ограничения:
\[
2k\le 11,\qquad k\le6.
\]
Так как $k$ целое, первое неравенство даёт $k\le5$, а второе — $k\le6$. Поэтому наибольшее допустимое значение — $5$.
\topic{Проверка}
При $k=5$ нужны десять двоек и пять троек; в $15!$ их соответственно одиннадцать и шесть — делимость выполняется. При $k=6$ потребовалось бы двенадцать двоек, но в $15!$ их только одиннадцать. Значит, увеличить $k$ нельзя.
\topic{Итог}
\finalword{} Числовой ответ $\boxed{k=5}$ правильный. Следует аккуратно отделять частичную информацию о простых множителях от полного разложения числа.
\topic{Задача на закрепление}
Найдите наибольшее целое неотрицательное $k$, при котором $12^k\mid16!$.
\newpage
\phantomsection\label{task:9}
\begin{center}{\LARGE\bfseries\color{deep} Задание № 9}\end{center}
\topic{Текст задания}
Найдите наибольшее целое неотрицательное $k$, при котором $18^k\mid20!$.
\topic{Ваше решение}
Правильно вычислены
\[
\vp{2}{20!}=10+5+2+1=18,\qquad\vp{3}{20!}=6+2=8,
\]
а также $\vp{5}{20!}=4$ и $\vp{7}{20!}=2$. Указано разложение
\[
20!=2^{18}\cdot3^8\cdot5^4\cdot7^2,
\]
а для степени основания записано $18^k=2^1\cdot3^2$. В конце получено $\min\{18/1,8/2\}=4$.
\topic{Ваш ответ}
$k=4$ — верно.
\topic{Ошибка}
\err{} В двух равенствах опущены необходимые простые множители или действие показателя степени. Эти записи в буквальном виде неверны.
\topic{Где именно возникает ошибка}
\textbf{Последний правильный шаг:} подсчёт показателей простых чисел в $20!$, в частности $\vp{3}{20!}=8$ и $\vp{7}{20!}=2$.

\textbf{Первый неправильный шаг:} запись $20!=2^{18}3^85^47^2$, где отсутствуют множители $11$, $13$, $17$, $19$. Позднее появляется ещё одно неверное равенство $18^k=2\cdot3^2$: правая часть равна $18$ при любом $k$.
\topic{Почему это неверно}
Числа $11$, $13$, $17$ и $19$ входят в произведение от $1$ до $20$, поэтому они обязательно входят в полное разложение $20!$. Кроме того, правило степени произведения имеет вид $(ab)^k=a^kb^k$. Если оставить показатели $1$ и $2$ без умножения на $k$, запись описывает число $18$, а не произвольную его степень. Например, при $k=2$ имеем $18^2=324=2^2\cdot3^4$, тогда как $2^1\cdot3^2=18$.
\topic{Ключевая идея}
\keyidea{} Сначала разложите \emph{основание} на простые множители, затем умножьте каждый его показатель на $k$. Для проверки делимости сравнивайте полученные показатели с показателями тех же простых чисел в факториале.
\topic{Исправление}
Два равенства нужно заменить на
\[
18^k=(2\cdot3^2)^k=2^k\cdot3^{2k},
\]
\[
20!=2^{18}\cdot3^8\cdot5^4\cdot7^2\cdot11\cdot13\cdot17\cdot19.
\]
Для определения $k$ достаточно первых двух показателей полного разложения.
\topic{Полное правильное решение}
По формуле Лежандра количество двоек в $20!$ равно
\[
\vp{2}{20!}=\floor{20/2}+\floor{20/4}+\floor{20/8}+\floor{20/16}
=10+5+2+1=18.
\]
Количество троек:
\[
\vp{3}{20!}=\floor{20/3}+\floor{20/9}=6+2=8.
\]
Так как $18^k=2^k3^{2k}$, необходимы и достаточны условия
\[
k\le18,\qquad 2k\le8.
\]
Второе условие даёт $k\le4$ и является более строгим. Следовательно, $k_{\max}=4$.
\topic{Проверка}
При $k=4$ число $18^4$ содержит четыре двойки и восемь троек; в $20!$ их восемнадцать и восемь. Делимость есть. При $k=5$ число $18^5$ содержит уже десять троек, что больше восьми; делимость пропадает.
\topic{Итог}
\finalword{} Правильное значение $\boxed{k=4}$ получено, однако два неверных промежуточных равенства важно исправить, чтобы обоснование стало строгим.
\topic{Задача на закрепление}
Найдите наибольшее целое неотрицательное $k$, при котором $50^k\mid25!$.
\newpage
\phantomsection\label{task:10}
\begin{center}{\LARGE\bfseries\color{deep} Задание № 10}\end{center}
\topic{Текст задания}
Найдите наибольшее целое неотрицательное $k$, при котором $72^k\mid30!$.
\topic{Ваше решение}
В начале записано $72^k=2^3\cdot3^2$. Затем верно вычислены показатели
\[
\vp{2}{30!}=15+7+3+1=26,\qquad\vp{3}{30!}=10+3+1=14.
\]
Для пятёрки записан промежуточный результат $7+1=8$ при сумме $\floor{30/5}+\floor{30/25}$. В разложении $30!$ выписаны только множители с основаниями $2$, $3$, $5$, $7$. Итоговый ответ — $7$.
\topic{Ваш ответ}
$k=7$ — верно.
\topic{Ошибка}
\err{} При возведении числа $72$ в степень не распределён показатель $k$; отдельно допущена арифметическая ошибка при подсчёте пятёрок в $30!$; частичное разложение факториала записано как полное.
\topic{Где именно возникает ошибка}
\textbf{Первый неправильный шаг:} $72^k=2^3\cdot3^2$. Справедливо только $72=2^3\cdot3^2$; для $72^k$ оба показателя должны содержать $k$.

\textbf{Следующий неправильный шаг:} $\floor{30/5}+\floor{30/25}=7+1=8$. На самом деле $\floor{30/5}=6$, поэтому результат равен $6+1=7$.

\textbf{Сохранённая правильная часть:} вычисления $\vp{2}{30!}=26$ и $\vp{3}{30!}=14$, а также итоговое ограничение для $k$.
\topic{Почему это неверно}
Правило степени произведения требует применять показатель $k$ к \emph{каждому} простому множителю: $(2^3\cdot3^2)^k=2^{3k}\cdot3^{2k}$. Формула Лежандра использует целую часть от деления. В последовательности $1,2,\ldots,30$ ровно шесть чисел, кратных пяти, и одно дополнительное кратное $25$. Следовательно, пятёрок в $30!$ семь. Более того, в полном разложении $30!$ присутствуют и простые числа больше семи. Это не влияет на максимальное $k$, поскольку $72$ содержит только простые множители $2$ и $3$.
\topic{Ключевая идея}
\keyidea{} Для делимости $a^k\mid n!$ разложите $a$ на простые множители и сравните требуемые показатели с соответствующими показателями в $n!$. Все вспомогательные вычисления должны оставаться точными, даже если не определяют итоговый ответ.
\topic{Исправление}
\[
72^k=(2^3\cdot3^2)^k=2^{3k}\cdot3^{2k},\qquad
\vp{5}{30!}=\floor{30/5}+\floor{30/25}=6+1=7.
\]
Если нужно полное разложение, верна запись
\[
\begin{aligned}
30!={}&2^{26}\cdot3^{14}\cdot5^7\cdot7^4\cdot11^2\cdot13^2\\
&\cdot17\cdot19\cdot23\cdot29.
\end{aligned}
\]
\topic{Полное правильное решение}
По формуле Лежандра
\[
\begin{aligned}
\vp{2}{30!}&=\floor{30/2}+\floor{30/4}+\floor{30/8}+\floor{30/16}
=15+7+3+1=26,\\
\vp{3}{30!}&=\floor{30/3}+\floor{30/9}+\floor{30/27}
=10+3+1=14.
\end{aligned}
\]
Для $72^k=2^{3k}3^{2k}$ получаем условия
\[
3k\le26,\qquad 2k\le14.
\]
Из первого $k\le\floor{26/3}=8$, из второго $k\le\floor{14/2}=7$. Совместно получаем наибольшее допустимое значение $k=7$.
\topic{Проверка}
При $k=7$ требуется $21$ двойка и $14$ троек; в $30!$ имеется $26$ двоек и $14$ троек, поэтому делимость выполняется. При $k=8$ потребовалось бы $16$ троек, а их только $14$. Следовательно, $k=8$ уже не подходит.
\topic{Итог}
\finalword{} Ответ $\boxed{k=7}$ верен. В дальнейшем особенно внимательно записывайте степень произведения и считайте целые части дробей.
\topic{Задача на закрепление}
Найдите наибольшее целое неотрицательное $k$, при котором $360^k\mid35!$.
\end{document}
